Learning-based controllers (like Reinforcement Learning) or human
operators are excellent at exploring complex tasks but often lack
formal safety guarantees.
\\
The idea of \textbf{Safety Filters}
is to filter the nominal input $u^\star$
and produce a safe input $u$
that minimally deviates from $u^\star$
while ensuring the system state remains in a safe set.

\subsection{Decoupling}

Instead of building safety directly into a complex objective function,
we decouple the problem.

\begin{enumerate}
  \item \textbf{Performance}
    A controller (e.g., a Learning-based Controller, human or AI input,...)
    focuses on the complex objectives and provides a suggested input $u_L(k)$.
  \item \textbf{Safety Filter} A verification module that sits between the
    controller and the system. It verifies $u_L(k)$ and modifies it
    \textbf{minimally} only if required to ensure constraint satisfaction
    for all future times.
\end{enumerate}

\subsection{Safe Set}

A central concept is the \textbf{Safe Set} $\mathcal{S}$
(usually Maximal Control Invariant Set),
which ensures states remain feasible indefinitely under
\textbf{Safe Control Law} $\pi_S(x)$.

\subsection{Implementation Types}

\subsubsection{Switching SF}

Binary choice. If $u_L$ leads to a safe state, use it.
Otherwise, \textit{fall back} to the safe input $Kx$.
This can lead to \textit{harsh} interventions.

\subsubsection{Minimally Invasive SF}

Formulated as a projection problem:
\[\min_{u_0} \|u_0 - u_L(k)\| \quad \text{s.t.} \quad f(x(k), u_0)
\in \mathcal{S}, \quad u_0 \in \mathcal{U}\]

\subsubsection{Control Barrier Function (CBF) SF}

Introduces a \textit{dampening} constraint to prevent
the system from rushing toward the boundary of the safe set.
It uses a parameter $\gamma \in (0,1]$ to reduce the
\textit{remaining budget} until the boundary:
\[V(f(x, u_0)) - V(x) \le \gamma(1 - V(x))\]
where $V$ is a CBF (e.g., derived from a Lyapunov function).

\subsection{Model Predictive Safety Filter (MPSF)}

While invariance-based filters are often restricted to small
ellipsoidal sets (which are conservative),
the \textbf{Predictive Safety Filter} uses MPC techniques to look further ahead.

\textbf{Mechanism}
It plans a safe forward trajectory of $N$ steps
that terminates in a safe set $\mathcal{X}_f$
(ensuring recursive feasibility and invariance).

\textbf{Implicit Safe Set:} The set of all states for which the MPSF
is feasible, $\mathcal{X}_{feas}$, is itself a control invariant set.
Crucially: $\mathcal{X}_f \subseteq \mathcal{X}_{feas} \subseteq
\mathcal{C}$. MPSF allows exploring a much larger state space than
$\mathcal{X}_f$.

\textbf{The Formulation}
\[\min_{u_0, \dots, u_{N-1}} \|u_0 - u_L(k)\|\]
\[\text{s.t. } x_{i+1} = f(x_i, u_i), \quad (x_i, u_i) \in
\mathcal{X} \times \mathcal{U}, \quad x_N \in \mathcal{X}_f\]
\textbf{Benefit}
Any state that is \textit{feasible} for this optimization
is considered safe.
This implicitly defines a much larger safe set
$\mathcal{X}_{feas}$ than a simple ellipsoidal $\mathcal{S}$.

\subsubsection{Robustness and Performance}

\textbf{Uncertainty}
In real-world applications
(like quadrotor control or autonomous racing),
safety depends on a perfect model.
To handle disturbances, the MPSF can be extended using
\textbf{Robust Tube MPC} (e.g., for bounded disturbances) or
\textbf{Stochastic MPC} techniques
(e.g., probabilistic constraint tightening).

\textbf{Softening the Intervention}
To prevent the filter from \textit{fighting} the learner,
a barrier term can be added to the learner's own cost function
to make it naturally prefer safe regions:
\[C_s(x,u) = C(x,u) - \gamma \log(1 - V(x))\]

\subsubsection{Advantages and Limitations}
Invariance-based SFs scale to high dimensions and are computationally
cheap but conservative.
MPSFs reduce conservatism (larger safe sets) but require solving an
MPC-like optimization online.

